\casechapter{Every Regular Problem Is Locally Gaussian}{Le Cam's theory of experiments}{Localize to a limit object}{ch:lecam}
\begin{refsection}[chapters/13-lecam.bib]

\newcommand{\lecamto}{\rightsquigarrow}
\newcommand{\lecamN}{\mathrm{N}}

\begin{chaptersummary}
Fisher claimed that no estimator can do better, asymptotically, than the maximum likelihood estimator. Hodges's estimator does better at a single point, and for two decades the claim survived only in restricted forms. Le Cam and Hájek dissolved the difficulty by zooming in: at the scale $\theta_0+h/\sqrt n$ every regular model looks like one Gaussian observation with mean $h$ and variance $1/I(\theta_0)$, where optimality questions are first exercises, and a transfer theorem carries the answers back. Superefficiency becomes the remark that a constant estimator is perfect at one point, and Fisher's bound becomes a local minimax theorem valid for every estimator.
\end{chaptersummary}

\section{Before: Fisher's bound and Hodges's estimator}\label{sec:lecam:before}

Let $X_1,\dots,X_n$ be independent observations from a distribution $P_\theta$ with density $p_\theta$ with respect to a $\sigma$-finite measure $\mu$, where $\theta$ ranges over an open interval $\Theta$. We write $P_\theta^n$ and $\E_\theta$ for the law of the sample and expectation under it; an \term{estimator} $T_n$ is a measurable function of the sample. When $\theta\mapsto\log p_\theta(x)$ is differentiable, its derivative $\dot\ell_\theta(x)$ is the \term{score}, and $I(\theta)=\E_\theta\dot\ell_\theta(X_1)^2$ is the \term{Fisher information}\index{Fisher information} in one observation. Throughout, $\Phi$ and $\varphi$ are the standard normal distribution function and density, and $Y_n\lecamto L$ means that the law of $Y_n$ converges weakly to $L$.

In two papers, \textcite{Fisher1922,Fisher1925} laid out a program for the maximum likelihood estimator $\hat\theta_n$ that, in the paraphrase of \textcite{Vovk2009}, makes two claims. First, $\sqrt n(\hat\theta_n-\theta)\lecamto\lecamN(0,1/I(\theta))$ under $P_\theta^n$. Second, if another estimator satisfies $\sqrt n(T_n-\theta)\lecamto\lecamN(0,v(\theta))$ under every $P^n_\theta$, then $v(\theta)\ge1/I(\theta)$. The first claim is a theorem, under hypotheses; the second, read literally, is false.

\textcite{Cramer1946} made the first claim rigorous, under hypotheses worth recording because they are what Le Cam later removed. With $\theta_0$ the true value, they are these: the support of $p_\theta$ does not depend on $\theta$; $\log p_\theta(x)$ has three continuous $\theta$-derivatives; $\int p_\theta\dd\mu=1$ may be differentiated twice under the integral sign; the third derivative is bounded near $\theta_0$ by a function $M$ with $\E_{\theta_0}M<\infty$; and $0<I(\theta_0)<\infty$. The proof finds a consistent root of the likelihood equation (not necessarily the maximizer), expands the equation at the root, applies the central limit theorem to the score and the law of large numbers to its derivative, and uses $M$ to kill the remainder. The conclusion depends on the model only through $I(\theta_0)$, and it says nothing about other estimators.

For the second claim the tool at hand was a finite-sample inequality, found independently by several authors in the 1940s, among them \textcite{Rao1945} and \textcite{Cramer1946}.\index{Cram\'er--Rao inequality}

\begin{theorem}[Cramér--Rao inequality]\label{thm:lecam:cramer-rao}
Suppose that $\theta\mapsto p_\theta(x)$ is differentiable for every $x$, that $0<I(\theta)<\infty$, and that differentiation in $\theta$ under the integral sign is allowed in $\int p_\theta\dd\mu=1$ and in $\E_\theta T=\int T\prod_ip_\theta(x_i)\dd\mu^n(x)$, where $T$ is an estimator with $\E_\theta T^2<\infty$. Then its bias $b(\theta)=\E_\theta T-\theta$ satisfies
\[
  \Var_\theta(T)\ \ge\ \frac{(1+b'(\theta))^2}{nI(\theta)}.
\]
In particular every unbiased estimator has variance at least $1/(nI(\theta))$.
\end{theorem}

\begin{proof}
Let $S=\sum_i\dot\ell_\theta(X_i)$ be the score of the sample. Differentiating $\int p_\theta\dd\mu=1$ shows that each $\dot\ell_\theta(X_i)$ has mean zero; these scores are independent with variance $I(\theta)$, so $\E_\theta S=0$ and $\Var_\theta(S)=nI(\theta)$. Differentiating $\E_\theta T$ gives $\E_\theta(TS)=1+b'(\theta)$, which is therefore $\Cov_\theta(T,S)$. The Cauchy--Schwarz inequality $\Cov_\theta(T,S)^2\le\Var_\theta(T)\Var_\theta(S)$ is the claim.
\end{proof}

The inequality has the wrong shape for the second claim. Unbiasedness constrains the mean of $T$ at every $\theta$, rarely holds for the maximum likelihood estimator, and is sometimes impossible: in the Bernoulli model the mean of every estimator is a polynomial in $p$ of degree at most $n$, so the log-odds has no unbiased estimator. The biased form involves $b'(\theta)$, a property of the estimator under judgment.

In 1951 Hodges produced a series of examples showing that the second claim fails; they first appeared in print in Le Cam's thesis \parencite[Section~1]{LeCam1953}. The version below became the standard one.\sidenote{The objection is older than the example. In a letter of 1930, quoted by \textcite{Vovk2009}, Hotelling pointed out to Fisher that a method of estimation tuned to the true value would have arbitrarily small variance there, and that nothing prevents its use when the true value is unknown. Hodges made the idea consistent at every parameter value.}

\begin{proposition}[Hodges's estimator]\label{prop:lecam:hodges}\index{Hodges's estimator}
Let $X_1,X_2,\dots$ be independent $\lecamN(\theta,1)$ variables, let $\bar X_n$ be the mean of the first $n$, and let $T_n=\bar X_n$ if $\abs{\bar X_n}\ge n^{-1/4}$ and $T_n=0$ otherwise. If $\theta\ne0$, then $\sqrt n(T_n-\theta)\lecamto\lecamN(0,1)$ and $n\E_\theta(T_n-\theta)^2\to1$. If $\theta=0$, then $\sqrt n\,T_n\to0$ in probability and $n\E_0T_n^2\to0$.
\end{proposition}

\begin{proof}
Under $P_\theta^n$ the variable $Z=\sqrt n(\bar X_n-\theta)$ is standard normal, and $T_n$ differs from $\bar X_n$ only on the event $A_n=\{\abs{\bar X_n}<n^{-1/4}\}$. Suppose first that $\theta\ne0$, and take $n$ so large that $n^{-1/4}\le\abs\theta/2$. On $A_n$ we then have $\abs{\bar X_n-\theta}>\abs\theta/2$, so $P_\theta^n(A_n)\le P(\abs Z>\sqrt n\abs\theta/2)\le2e^{-n\theta^2/8}$. Outside $A_n$ the variable $\sqrt n(T_n-\theta)$ equals $Z$, and on $A_n$ it equals $-\sqrt n\,\theta$. Hence $\sqrt n(T_n-\theta)\lecamto\lecamN(0,1)$, and
\[
  n\E_\theta(T_n-\theta)^2=1-\E[Z^2\ind_{A_n}]+n\theta^2P^n_\theta(A_n)\to1,
\]
because $\E[Z^2\ind_{A_n}]\le(3P^n_\theta(A_n))^{1/2}$ by the Cauchy--Schwarz inequality and $n\theta^2e^{-n\theta^2/8}\to0$. If $\theta=0$, then $\sqrt n\,T_n=Z\ind\{\abs Z\ge n^{1/4}\}$, which vanishes outside an event of probability $P(\abs Z\ge n^{1/4})\to0$, and $n\E_0T_n^2=\E[Z^2\ind\{\abs Z\ge n^{1/4}\}]\to0$ by dominated convergence.
\end{proof}

So $v(\theta)=1=1/I(\theta)$ for $\theta\ne0$, while $v(0)=0$. A parameter value at which an estimator beats Fisher's bound in this way is a point of \term{superefficiency}\index{superefficiency}; shrinking toward finitely many points produces finitely many of them. \Cref{thm:lecam:cramer-rao} does not forbid this, because the bias of $T_n$ has derivative tending to $-1$ at the origin (Exercise~\ref{ex:lecam:hodges}).

The responses, in the classification of \textcite{Vovk2009}, took three forms. The first proved the second claim for restricted classes of estimators, such as smooth functions of observed frequencies or estimators whose convergence to normality is uniform in $\theta$.\sidenote{Wolfowitz objected that proving the maximum likelihood estimator best by declaring strong competitors ineligible is like becoming a town's chess champion after the better players have been barred; his analogy is quoted by \textcite{Vovk2009}.} The second showed that the damage is small: Le Cam announced at the end of 1951, and treated in his thesis \parencite{LeCam1953}, that superefficiency can happen only on a Lebesgue-null set; the proof in the thesis was flawed, and correct proofs came later \parencite{Vovk2009,vanderVaart1997}. The third judged an estimator by its worst performance over shrinking neighborhoods of each point. Le Cam's thesis has results of this kind too, among them that a superefficient estimator performs poorly near its points of superefficiency, and the line culminated in Hájek's local asymptotic minimax theorem \parencite{Hajek1972}.\sidenote{Lucien Le Cam (1924--2000) was born in Croze, in the Creuse. His thesis, the published form of his 1952 Berkeley dissertation, also proves local asymptotic minimax and admissibility properties of Bayes estimates and transfers them to the maximum likelihood estimator.} Here is the shape of a proof of the first kind, built from tools available in 1950.

\begin{aside}[A restricted-class proof of Fisher's bound]
Fix $\theta_0$ and $h>0$, put $\theta_n=\theta_0+h/\sqrt n$, and \emph{assume} that $P^n_{\vartheta_n}(\sqrt n(T_n-\vartheta_n)\le x)\to\Phi(x/\sqrt v)$ for every $x$ and every sequence $\vartheta_n$ with $\sqrt n(\vartheta_n-\theta_0)$ bounded; Hodges's estimator violates this at $\theta_0=0$. The test that rejects $\theta_0$ when $T_n\ge\theta_n$ has level tending to $\alpha=1-\Phi(h/\sqrt v)$ and power tending to $\tfrac12$, so by the Neyman--Pearson lemma the likelihood ratio test at the same level has asymptotic power at least $\tfrac12$. To compute that power, expand $\Lambda_n=\sum_i\log(p_{\theta_n}/p_{\theta_0})(X_i)$ to second order with a third-derivative remainder bound as in Cramér's theorem, and apply the central limit theorem under $\theta_0$ and, for triangular arrays, under $\theta_n$. The outcome, $\Lambda_n\approx\lecamN(\mp\tfrac12h^2I,h^2I)$ under $\theta_0$ and $\theta_n$ with $I=I(\theta_0)$, gives the power $1-\Phi(h/\sqrt v-h\sqrt I)$, and comparing with $\tfrac12$ yields $v\ge1/I$. The argument covers only estimators that satisfy the assumption, and its expansions must be redone in each model.
\end{aside}

By the late 1950s, then, one had a normality theorem under heavy assumptions, an inequality blind to Hodges's example, a null-set theorem, and bounds for classes restricted to make proofs work. Questions remained. In what sense, with no restriction on the competitors, is the maximum likelihood estimator optimal? Why does the answer depend on the model only through $I(\theta_0)$? Is there one theory for all losses, for tests, and for confidence intervals? And which of Cramér's hypotheses are needed at all?

\section{The move: zoom in to the scale of the noise}\label{sec:lecam:move}

Hodges's estimator is excellent at the origin and, as we shall see, poor at the points $h/\sqrt n$ nearby. Every statement about a single parameter value is exposed to this trick. So we study a neighborhood instead of a point, choose its size, and find what the problem looks like at that size.

Fix $\theta_0\in\Theta$ and write the parameter as $\theta=\theta_0+h/\sqrt n$. We call $h$ the \term{local parameter}\index{localization} and the family $\cE_n=(P^n_{\theta_0+h/\sqrt n}:h\in H_n)$, with $H_n=\{h:\theta_0+h/\sqrt n\in\Theta\}$ increasing to $\R$, the \term{local experiment} at $\theta_0$. An estimator $T_n$ of $\theta$ becomes the estimator $\sqrt n(T_n-\theta_0)$ of $h$, and the normalized loss $n(T_n-\theta)^2$ is its squared error.\sidenote{In Le Cam's language an \emph{experiment} is a family of probability distributions on a common space, indexed by a parameter. His deficiency distance makes the experiments with a given parameter set, up to equivalence, into a metric space \parencite{LeCam1986}.}

The scale $1/\sqrt n$ is forced. In the normal model the Hellinger affinity $\int\sqrt{pq}\dd\mu$ of $\lecamN(\theta,1)^n$ and $\lecamN(\theta',1)^n$ is $\exp(-n(\theta-\theta')^2/8)$, and the smallest sum of the two error probabilities of a test of $\theta$ against $\theta'$ tends to zero if $\sqrt n\abs{\theta-\theta'}\to\infty$ and to one if $\sqrt n\abs{\theta-\theta'}\to0$ (Exercise~\ref{ex:lecam:hellinger}). At a coarser scale the local problem becomes trivial and at a finer one hopeless; at the scale $1/\sqrt n$ it neither trivializes nor degenerates, and the same holds in every model that is regular in the sense of \cref{sec:lecam:limit}.

What does the local experiment look like as $n\to\infty$? The answer is the humblest object in statistics.

\begin{definition}[Gaussian shift]\label{def:lecam:gaussian-shift}
For $\sigma>0$, the \term{Gaussian shift experiment}\index{Gaussian shift experiment} is the family $(\lecamN(h,\sigma^2):h\in\R)$: one observation $X$ from a normal distribution with unknown mean $h$ and known variance $\sigma^2$.
\end{definition}

This is the first estimation problem of every statistics course, and it is where we stoop low. Its log-likelihood ratio
\[
  \log\frac{\dd\lecamN(h,\sigma^2)}{\dd\lecamN(0,\sigma^2)}(X)=h\,\frac X{\sigma^2}-\frac{h^2}{2\sigma^2}
\]
is a quadratic function of $h$ whose linear coefficient is normal with mean zero and variance $1/\sigma^2$ under $h=0$. The next definition asks the local experiment to have, asymptotically, log-likelihood ratios of exactly this form.

\begin{definition}[local asymptotic normality]\label{def:lecam:lan}
The model is \term{locally asymptotically normal}\index{local asymptotic normality} (LAN) at $\theta_0$ with information $I_0>0$ if there are statistics $\Delta_n=\Delta_n(X_1,\dots,X_n)$ with $\Delta_n\lecamto\lecamN(0,I_0)$ under $P^n_{\theta_0}$ such that, for every $h\in\R$,
\begin{equation}\label{eq:lecam:lan}
  \log\prod_{i=1}^n\frac{p_{\theta_0+h/\sqrt n}(X_i)}{p_{\theta_0}(X_i)}=h\Delta_n-\tfrac12h^2I_0+o_{P}(1)\qquad\text{under }P^n_{\theta_0}.
\end{equation}
\end{definition}

With $\sigma^2=1/I_0$ the Gaussian shift has log-likelihood ratio $h\,(I_0X)-\frac12h^2I_0$, with $I_0X\sim\lecamN(0,I_0)$ under $h=0$. So \eqref{eq:lecam:lan} says that $\Delta_n$ plays the role of $I_0X$, and $\Delta_n/I_0$ is, asymptotically, the observation of the Gaussian shift $(\lecamN(h,1/I_0):h\in\R)$. \textcite{LeCam1960} introduced the notion. In the normal model $\lecamN(\theta,1)$ the expansion is exact, with $\Delta_n=\sqrt n(\bar X_n-\theta_0)$, $I_0=1$, and no remainder.

\begin{definition}[convergence of experiments]\label{def:lecam:convergence}
Experiments $(P_{n,h}:h\in H_n)$, with $H_n$ increasing to $H$, \term{converge}\index{convergence of experiments} to an experiment $(Q_h:h\in H)$, the \term{limit experiment}\index{limit experiment}, if for every finite set $F\subset H$ and every $h_0\in H$ the vector $(\dd P_{n,h}/\dd P_{n,h_0})_{h\in F}$ under $P_{n,h_0}$ converges in distribution to $(\dd Q_h/\dd Q_{h_0})_{h\in F}$ under $Q_{h_0}$.
\end{definition}

Le Cam's original definition uses the deficiency, a distance measuring how well the risk functions of one experiment can be matched in another; for finite parameter sets, convergence in this distance can be expressed through the laws of likelihood ratios \parencite{LeCam1986,vanderVaart1998}. Under LAN the local experiments converge to the Gaussian shift with $\sigma^2=1/I_0$; at the base point $0$ this is \eqref{eq:lecam:lan}, and contiguity (\cref{sec:lecam:transfer}) handles the others.

The move rests on one principle, made precise in \cref{thm:lecam:representation}: \emph{whatever a sequence of estimators achieves asymptotically along local alternatives, some randomized estimator achieves in the limit experiment.} Lower bounds proved in the limit therefore bind all estimator sequences, and optimal procedures in the limit are targets for sequences.

What would one have had to notice? Hodges's estimator and Le Cam's thesis showed that performance must be judged over shrinking neighborhoods, and the Hellinger computation fixes their size. And the arguments of \cref{sec:lecam:before} computed the same quadratic $h\Delta_n-\frac12h^2I$, with $\Delta_n$ a normalized sum of scores, and threw it away after one use. The move keeps the quadratic, recognizes it as the log-likelihood of one Gaussian observation, and makes the Gaussian shift the object of study.

\section{After: the Gaussian shift does the work}\label{sec:lecam:limit}

\subsection{Hodges's estimator, seen from nearby}

\begin{proposition}[local risk of Hodges's estimator]\label{prop:lecam:hodges-local}
In the setting of \cref{prop:lecam:hodges}, fix $h\in\R$, put $\theta_n=h/\sqrt n$, and let $R_n(h)=n\E_{\theta_n}(T_n-\theta_n)^2$. Then $\sqrt n(T_n-\theta_n)\to-h$ in $P^n_{\theta_n}$-probability, and $R_n(h)\to h^2$. Moreover, $\sup_{\theta\in\R}n\E_\theta(T_n-\theta)^2>\sqrt n/3$ for every $n$.
\end{proposition}

\begin{proof}
Under $P^n_{\theta_n}$ the variable $Y=\sqrt n\,\bar X_n$ has the law $\lecamN(h,1)$, which does not depend on $n$, and
\[
  \sqrt n(T_n-\theta_n)=(Y-h)\ind\{\abs Y\ge n^{1/4}\}-h\ind\{\abs Y<n^{1/4}\}.
\]
Since $P(\abs Y\ge n^{1/4})\to0$, the first claim follows. The two indicators are disjoint, so
\[
  R_n(h)=\E\big[(Y-h)^2\ind\{\abs Y\ge n^{1/4}\}\big]+h^2P(\abs Y<n^{1/4});
\]
the first term tends to zero by dominated convergence and the second to $h^2$. For the last claim take $h=n^{1/4}$, that is, $\theta=n^{-1/4}$. The second term alone equals $\sqrt n\,P(-2n^{1/4}<Y-h<0)$, which is at least $\sqrt n\,(\tfrac12-\Phi(-2))>\sqrt n/3$.
\end{proof}

In the local experiment, then, the Hodges sequence converges to the constant estimator $T=0$ of the Gaussian shift $(\lecamN(h,1):h\in\R)$, whose risk $h^2$ is zero at $h=0$, below the risk $1$ of $X$ when $\abs h<1$, and above it when $\abs h>1$. Superefficiency is nothing more than this: a constant is a perfect estimator at one point and pays everywhere else. \Cref{fig:lecam:hodges} shows the bill, a worst local risk over $\abs h\le c$ that tends to $c^2$ and a worst global risk that grows like $\sqrt n$. The fixed-$\theta$ analysis of \cref{prop:lecam:hodges} saw none of this, because each fixed $\theta\ne0$ eventually leaves the window of width about $n^{-1/4}$ in which the estimator misbehaves.

\begin{figure}[t]
\centering
\begin{tikzpicture}[x=0.86cm,y=0.36cm]
  \draw[muted,->] (-6.4,0) -- (6.6,0) node[right,font=\small,text=ink] {$h$};
  \draw[muted,->] (0,0) -- (0,10.4) node[above,font=\small,text=ink] {local risk};
  \foreach \x in {-6,-4,-2,2,4,6} \draw[muted] (\x,0) -- (\x,-0.25) node[below,font=\scriptsize,text=ink] {$\x$};
  \foreach \y in {1,4,8} \draw[muted] (0,\y) -- (-0.12,\y) node[left,font=\scriptsize,text=ink] {$\y$};
  \draw[sandline,line width=0.9pt] (-6,1) -- (6,1);
  \node[font=\scriptsize,text=muted,anchor=north east] at (6.0,0.92) {sample mean};
  \draw[muted,dashed,domain=-3.2:3.2,samples=80] plot (\x,{\x*\x});
  \node[font=\scriptsize,text=muted,anchor=east] at (-3.25,10.0) {$h^2$};
  \draw[sea!45,line width=0.8pt] plot coordinates {(-6,1.00) (-5.9,1.00) (-5.8,1.00) (-5.7,1.00) (-5.6,1.00) (-5.5,1.00) (-5.4,1.01) (-5.3,1.01) (-5.2,1.01) (-5.1,1.01) (-5,1.02) (-4.9,1.03) (-4.8,1.03) (-4.7,1.04) (-4.6,1.06) (-4.5,1.08) (-4.4,1.10) (-4.3,1.12) (-4.2,1.15) (-4.1,1.19) (-4,1.23) (-3.9,1.28) (-3.8,1.34) (-3.7,1.41) (-3.6,1.48) (-3.5,1.56) (-3.4,1.64) (-3.3,1.73) (-3.2,1.83) (-3.1,1.93) (-3,2.03) (-2.9,2.12) (-2.8,2.22) (-2.7,2.30) (-2.6,2.38) (-2.5,2.44) (-2.4,2.49) (-2.3,2.52) (-2.2,2.54) (-2.1,2.53) (-2,2.50) (-1.9,2.45) (-1.8,2.38) (-1.7,2.28) (-1.6,2.17) (-1.5,2.04) (-1.4,1.90) (-1.3,1.75) (-1.2,1.59) (-1.1,1.42) (-1,1.26) (-0.9,1.09) (-0.8,0.94) (-0.7,0.79) (-0.6,0.66) (-0.5,0.54) (-0.4,0.44) (-0.3,0.37) (-0.2,0.31) (-0.1,0.27) (0,0.26) (0.1,0.27) (0.2,0.31) (0.3,0.37) (0.4,0.44) (0.5,0.54) (0.6,0.66) (0.7,0.79) (0.8,0.94) (0.9,1.09) (1,1.26) (1.1,1.42) (1.2,1.59) (1.3,1.75) (1.4,1.90) (1.5,2.04) (1.6,2.17) (1.7,2.28) (1.8,2.38) (1.9,2.45) (2,2.50) (2.1,2.53) (2.2,2.54) (2.3,2.52) (2.4,2.49) (2.5,2.44) (2.6,2.38) (2.7,2.30) (2.8,2.22) (2.9,2.12) (3,2.03) (3.1,1.93) (3.2,1.83) (3.3,1.73) (3.4,1.64) (3.5,1.56) (3.6,1.48) (3.7,1.41) (3.8,1.34) (3.9,1.28) (4,1.23) (4.1,1.19) (4.2,1.15) (4.3,1.12) (4.4,1.10) (4.5,1.08) (4.6,1.06) (4.7,1.04) (4.8,1.03) (4.9,1.03) (5,1.02) (5.1,1.01) (5.2,1.01) (5.3,1.01) (5.4,1.01) (5.5,1.00) (5.6,1.00) (5.7,1.00) (5.8,1.00) (5.9,1.00) (6,1.00)};
  \draw[sea,line width=0.8pt] plot coordinates {(-6,1.03) (-5.9,1.05) (-5.8,1.06) (-5.7,1.08) (-5.6,1.11) (-5.5,1.14) (-5.4,1.18) (-5.3,1.23) (-5.2,1.28) (-5.1,1.35) (-5,1.44) (-4.9,1.54) (-4.8,1.65) (-4.7,1.78) (-4.6,1.93) (-4.5,2.09) (-4.4,2.27) (-4.3,2.47) (-4.2,2.68) (-4.1,2.91) (-4,3.14) (-3.9,3.38) (-3.8,3.62) (-3.7,3.85) (-3.6,4.08) (-3.5,4.30) (-3.4,4.49) (-3.3,4.66) (-3.2,4.81) (-3.1,4.92) (-3,5.00) (-2.9,5.04) (-2.8,5.04) (-2.7,5.00) (-2.6,4.92) (-2.5,4.81) (-2.4,4.65) (-2.3,4.47) (-2.2,4.26) (-2.1,4.02) (-2,3.77) (-1.9,3.50) (-1.8,3.22) (-1.7,2.93) (-1.6,2.64) (-1.5,2.36) (-1.4,2.08) (-1.3,1.82) (-1.2,1.57) (-1.1,1.33) (-1,1.11) (-0.9,0.91) (-0.8,0.72) (-0.7,0.56) (-0.6,0.42) (-0.5,0.30) (-0.4,0.20) (-0.3,0.13) (-0.2,0.07) (-0.1,0.04) (0,0.03) (0.1,0.04) (0.2,0.07) (0.3,0.13) (0.4,0.20) (0.5,0.30) (0.6,0.42) (0.7,0.56) (0.8,0.72) (0.9,0.91) (1,1.11) (1.1,1.33) (1.2,1.57) (1.3,1.82) (1.4,2.08) (1.5,2.36) (1.6,2.64) (1.7,2.93) (1.8,3.22) (1.9,3.50) (2,3.77) (2.1,4.02) (2.2,4.26) (2.3,4.47) (2.4,4.65) (2.5,4.81) (2.6,4.92) (2.7,5.00) (2.8,5.04) (2.9,5.04) (3,5.00) (3.1,4.92) (3.2,4.81) (3.3,4.66) (3.4,4.49) (3.5,4.30) (3.6,4.08) (3.7,3.85) (3.8,3.62) (3.9,3.38) (4,3.14) (4.1,2.91) (4.2,2.68) (4.3,2.47) (4.4,2.27) (4.5,2.09) (4.6,1.93) (4.7,1.78) (4.8,1.65) (4.9,1.54) (5,1.44) (5.1,1.35) (5.2,1.28) (5.3,1.23) (5.4,1.18) (5.5,1.14) (5.6,1.11) (5.7,1.08) (5.8,1.06) (5.9,1.05) (6,1.03)};
  \draw[seadeep,line width=0.8pt] plot coordinates {(-6,1.69) (-5.9,1.85) (-5.8,2.03) (-5.7,2.24) (-5.6,2.49) (-5.5,2.76) (-5.4,3.06) (-5.3,3.40) (-5.2,3.76) (-5.1,4.15) (-5,4.57) (-4.9,5.00) (-4.8,5.44) (-4.7,5.88) (-4.6,6.33) (-4.5,6.76) (-4.4,7.18) (-4.3,7.57) (-4.2,7.92) (-4.1,8.24) (-4,8.50) (-3.9,8.71) (-3.8,8.86) (-3.7,8.96) (-3.6,8.99) (-3.5,8.95) (-3.4,8.86) (-3.3,8.72) (-3.2,8.51) (-3.1,8.26) (-3,7.97) (-2.9,7.64) (-2.8,7.29) (-2.7,6.90) (-2.6,6.50) (-2.5,6.09) (-2.4,5.68) (-2.3,5.26) (-2.2,4.84) (-2.1,4.44) (-2,4.04) (-1.9,3.66) (-1.8,3.29) (-1.7,2.93) (-1.6,2.60) (-1.5,2.29) (-1.4,1.99) (-1.3,1.72) (-1.2,1.46) (-1.1,1.23) (-1,1.01) (-0.9,0.82) (-0.8,0.65) (-0.7,0.50) (-0.6,0.36) (-0.5,0.25) (-0.4,0.16) (-0.3,0.09) (-0.2,0.04) (-0.1,0.01) (0,0.00) (0.1,0.01) (0.2,0.04) (0.3,0.09) (0.4,0.16) (0.5,0.25) (0.6,0.36) (0.7,0.50) (0.8,0.65) (0.9,0.82) (1,1.01) (1.1,1.23) (1.2,1.46) (1.3,1.72) (1.4,1.99) (1.5,2.29) (1.6,2.60) (1.7,2.93) (1.8,3.29) (1.9,3.66) (2,4.04) (2.1,4.44) (2.2,4.84) (2.3,5.26) (2.4,5.68) (2.5,6.09) (2.6,6.50) (2.7,6.90) (2.8,7.29) (2.9,7.64) (3,7.97) (3.1,8.26) (3.2,8.51) (3.3,8.72) (3.4,8.86) (3.5,8.95) (3.6,8.99) (3.7,8.96) (3.8,8.86) (3.9,8.71) (4,8.50) (4.1,8.24) (4.2,7.92) (4.3,7.57) (4.4,7.18) (4.5,6.76) (4.6,6.33) (4.7,5.88) (4.8,5.44) (4.9,5.00) (5,4.57) (5.1,4.15) (5.2,3.76) (5.3,3.40) (5.4,3.06) (5.5,2.76) (5.6,2.49) (5.7,2.24) (5.8,2.03) (5.9,1.85) (6,1.69)};
  \node[font=\scriptsize,text=sea!80,anchor=west] at (2.45,2.75) {$n=16$};
  \node[font=\scriptsize,text=sea,anchor=west] at (3.1,5.25) {$n=81$};
  \node[font=\scriptsize,text=seadeep,anchor=west] at (3.95,9.0) {$n=256$};
\end{tikzpicture}
\caption{Local risk $h\mapsto n\E_{h/\sqrt n}(T_n-h/\sqrt n)^2$ of Hodges's estimator for $n=16$, $81$, and $256$ (light to dark), against the constant risk $1$ of the sample mean and the pointwise limit $h^2$ (dashed). The dip at $h=0$ is paid for by bumps of height about $\sqrt n/2$ near $\abs h=n^{1/4}$.}
\label{fig:lecam:hodges}
\end{figure}

\subsection{Minimax risk of the Gaussian shift}

Could some estimator beat the observation $X$ uniformly over a neighborhood? The answer takes one inequality, which we prove for a general family because \cref{sec:lecam:transfer} uses it again. It is a Bayesian Cramér--Rao bound, due to \textcite{vanTrees1968}: a prior replaces unbiasedness, and an integration by parts in $\theta$ replaces the differentiation of the bias.\sidenote{A prior under which every estimator does badly on average certifies a lower bound on the worst case. This is Yao's principle from \cref{ch:lpduality}, used at the scale of the noise.}

\begin{lemma}[van Trees inequality]\label{lem:lecam:vantrees}\index{van Trees inequality}
Let $(p_\theta:\theta\in[a,b])$ be probability densities with respect to $\mu$ such that
\begin{enumerate}[label=(\alph*)]
\item for every $x$ the map $\theta\mapsto p_\theta(x)$ is continuously differentiable on $[a,b]$, with derivative $\dot p_\theta(x)$;
\item for every $\theta$, $\int\dot p_\theta\dd\mu=0$, and the information $I(\theta)=\int(\dot p_\theta^2/p_\theta)\dd\mu$ is bounded on $[a,b]$.
\end{enumerate}
Let $\lambda$ be a probability density on $[a,b]$ that is absolutely continuous, vanishes at $a$ and $b$, and has finite information $J(\lambda)=\int_a^b(\lambda'^2/\lambda)\dd\theta$, with $0/0=0$ here and in $I$. Then every estimator $T$ satisfies
\[
  \int_a^b\E_\theta(T-\theta)^2\,\lambda(\theta)\dd\theta\ \ge\ \frac1{\int_a^bI(\theta)\lambda(\theta)\dd\theta+J(\lambda)}.
\]
\end{lemma}

\begin{proof}
The plan is to write $1$ as the integral of $T-\theta$ against $\partial_\theta(p_\theta\lambda)$ and to apply the Cauchy--Schwarz inequality. Put $f(x,\theta)=p_\theta(x)\lambda(\theta)$, and let $B$ be the left side, the Bayes risk of $T$. If $B=\infty$, then there is nothing to prove. For each $x$ the function $f(x,\cdot)$ is absolutely continuous, with derivative $\partial_\theta f=\dot p_\theta\lambda+p_\theta\lambda'$ almost everywhere, and it vanishes at $a$ and $b$ because $\lambda$ does. Integration by parts therefore gives, for every $x$,
\begin{equation}\label{eq:lecam:ibp}
  \int_a^b\big(T(x)-\theta\big)\,\partial_\theta f(x,\theta)\dd\theta=\int_a^bf(x,\theta)\dd\theta.
\end{equation}
Next we compute the information $C=\iint(\partial_\theta f)^2/f\dd\mu\dd\theta$ of $f$, with the integrand set to zero where $f=0$; there $\partial_\theta f=0$ almost everywhere, because a nonnegative absolutely continuous function has derivative zero at almost every one of its zeros. Expanding the square and integrating over $x$,
\[
  C=\int_a^bI\lambda\dd\theta+2\int_a^b\lambda'(\theta)\Big(\int\dot p_\theta\dd\mu\Big)\dd\theta+J(\lambda)=\int_a^bI\lambda\dd\theta+J(\lambda),
\]
because the inner integral vanishes by (b); the middle term is absolutely integrable because $\int\abs{\dot p_\theta}\dd\mu\le I(\theta)^{1/2}$, $I$ is bounded, and $\lambda'$ is integrable. By the Cauchy--Schwarz inequality, $\iint\abs{T-\theta}\abs{\partial_\theta f}\le(BC)^{1/2}<\infty$, so we may integrate \eqref{eq:lecam:ibp} over $x$ and exchange the order of integration. The right side integrates to $1$ and the left side is at most $(BC)^{1/2}$, which proves $B\ge1/C$.
\end{proof}

The lemma covers a randomized estimator $T=t(X,V)$, with $V$ uniform on $[0,1]$ and independent of $X$, through the densities $p_\theta(x)\ind_{[0,1]}(v)$.

\begin{theorem}[minimax risk of the Gaussian shift]\label{thm:lecam:gaussian-minimax}
Let $X\sim\lecamN(h,\sigma^2)$ with $\sigma>0$ known. For every $c>0$ and every estimator $T$, randomized or not,
\[
  \sup_{\abs h\le c}\E_h(T-h)^2\ \ge\ \frac{\sigma^2c^2}{c^2+\pi^2\sigma^2}.
\]
Consequently the minimax risk over $\abs h\le c$ tends to $\sigma^2$ as $c\to\infty$, and $X$ is minimax over $h\in\R$, with risk $\sigma^2$.
\end{theorem}

\begin{proof}
Apply \cref{lem:lecam:vantrees} on $[-c,c]$ to the densities of $\lecamN(h,\sigma^2)$, whose information is $1/\sigma^2$ for every $h$, with the prior $\lambda_c(h)=c^{-1}\cos^2(\pi h/(2c))$. This density vanishes at $\pm c$, and $\lambda_c'^2/\lambda_c=(\pi^2/c^3)\sin^2(\pi h/(2c))$, so $J(\lambda_c)=\pi^2/c^2$. The supremum of the risk over $[-c,c]$ is at least its $\lambda_c$-average, which the lemma bounds below by $1/(\sigma^{-2}+\pi^2/c^2)$. This bound tends to $\sigma^2$, the constant risk of $X$, as $c\to\infty$, and the supremum over $\R$ is at least the supremum over every $[-c,c]$.
\end{proof}

The cosine-squared prior is the best one for this argument: every density on $[-c,c]$ that vanishes at the endpoints has $J(\lambda)\ge\pi^2/c^2$ (Exercise~\ref{ex:lecam:prior}).\sidenote{The square root of $\lambda_c$ is the first Dirichlet eigenfunction of $-\mathrm d^2/\mathrm dh^2$ on $[-c,c]$, and $J(\lambda)=4\int\big((\sqrt\lambda)'\big)^2$, so the exercise is Wirtinger's inequality.}

\subsection{Equivariant estimators in the limit}

\Cref{thm:lecam:gaussian-minimax} compares worst cases. Estimators whose error has the same law wherever the truth is admit a finer comparison, and in the Gaussian shift it is complete. Write $\varphi_\sigma$ for the density of $\lecamN(0,\sigma^2)$.

\begin{proposition}[convolution in the Gaussian shift]\label{prop:lecam:convolution-limit}
Let $X\sim\lecamN(h,\sigma^2)$, let $V$ be uniform on $[0,1]$ and independent of $X$, and let $T=t(X,V)$ for a measurable function $t$. Suppose that the law of $T-h$ under $h$ is the same distribution $L$ for every $h\in\R$. Then $T-X$ is independent of $X$, and $L=\lecamN(0,\sigma^2)*M$, with $M$ the law of $T-X$.
\end{proposition}

\begin{proof}
The plan is to compare characteristic functions and to use the completeness of the normal location family. Fix $s\in\R$, let $\psi$ be the characteristic function of $L$, and put $g_s(x)=\E e^{is\,t(x,V)}$, a measurable function bounded by one. The hypothesis says that $\E_he^{isT}=e^{ish}\psi(s)$ for every $h$, while $\E_he^{isX}=e^{ish-\sigma^2s^2/2}$. Hence the bounded function $f_s(x)=g_s(x)-\psi(s)e^{\sigma^2s^2/2}e^{isx}$ satisfies $\int f_s(x)\varphi_\sigma(x-h)\dd x=0$ for every $h$. After division by $\varphi_\sigma(h)$, this says that the function $F(z)=\int f_s(x)e^{-x^2/2\sigma^2}e^{zx}\dd x$ vanishes on the real axis. It is entire, so it vanishes on the imaginary axis too, where it is the Fourier transform of the integrable function $f_s(x)e^{-x^2/2\sigma^2}$; by Fourier uniqueness, $f_s=0$ almost everywhere.

So for each $s$, for almost every $x$, the characteristic function of $t(x,V)-x$ at $s$ equals $\chi(s)=\psi(s)e^{\sigma^2s^2/2}$. Discarding a null set of $x$ for each rational $s$ and using continuity in $s$, we find a null set $N$ such that for $x\notin N$ the law of $t(x,V)-x$ has characteristic function $\chi$, and so is one and the same distribution $M$. Since $V$ is independent of $X$, that law is the conditional law of $T-X$ given $X=x$. Hence $T-X$ is independent of $X$ with law $M$, and $T-h=(X-h)+(T-X)$ has the law $\lecamN(0,\sigma^2)*M$.
\end{proof}

An equivariant estimator is thus the observation plus independent noise, and its risk is $\E_h(T-h)^2=\sigma^2+\int m^2\,M(\dd m)\ge\sigma^2$, with equality only for $T=X$; the same comparison holds for every loss that is symmetric and nondecreasing in $\abs{T-h}$ \parencite[Chapter~8]{vanderVaart1998}.\sidenote{Le Cam's own proof of Hájek's convolution theorem combined limits of experiments with a convolution result from C.~H. Boll's 1955 Stanford thesis \parencite{LeCam1972}.}

\subsection{Which models are locally Gaussian}

For exponential families, reaching the limit object takes a Taylor expansion.

\begin{theorem}[LAN for exponential families]\label{thm:lecam:lan-expfam}\index{exponential family}
Let $p_\theta(x)=\exp(\theta S(x)-A(\theta))$ with respect to $\mu$, for $\theta$ in an open interval $\Theta$ on which $A(\theta)=\log\int e^{\theta S}\dd\mu$ is finite, and suppose that $S(X_1)$ is not almost surely constant. Then the model is LAN at every $\theta_0\in\Theta$, with $\Delta_n=n^{-1/2}\sum_i(S(X_i)-A'(\theta_0))$ and $I_0=A''(\theta_0)=I(\theta_0)$. Moreover, the remainder in \eqref{eq:lecam:lan} is a number $r_n(h)$, not a random variable, and $\sup_{\abs h\le c}\abs{r_n(h)}\to0$ for every $c>0$.
\end{theorem}

\begin{proof}
A standard dominated-convergence argument shows that $A$ is infinitely differentiable on $\Theta$, with $A'(\theta)=\E_\theta S$ and $A''(\theta)=\Var_\theta S>0$. The score is $\dot\ell_\theta=S-A'(\theta)$, so $I(\theta)=A''(\theta)$. The logarithm on the left equals $(h/\sqrt n)\sum_iS(X_i)-n[A(\theta_0+h/\sqrt n)-A(\theta_0)]$ exactly, and adding and subtracting $\sqrt n\,hA'(\theta_0)$ rewrites it as
\[
  h\Delta_n-n\Big[A\big(\theta_0+\tfrac h{\sqrt n}\big)-A(\theta_0)-A'(\theta_0)\tfrac h{\sqrt n}\Big].
\]
By Taylor's theorem the bracket equals $A''(\xi)h^2/(2n)$ for some $\xi$ between $\theta_0$ and $\theta_0+h/\sqrt n$, so $r_n(h)=\tfrac12h^2(A''(\theta_0)-A''(\xi))$, which tends to zero uniformly in $\abs h\le c$ because $A''$ is continuous. Finally, $\Delta_n\lecamto\lecamN(0,A''(\theta_0))$ under $\theta_0$ by the central limit theorem.
\end{proof}

Compare Cramér's route in \cref{sec:lecam:before}: no consistency argument and no third derivative are needed, because the expansion is exact up to a deterministic remainder.

For the normal location model ($S(x)=x$, $\mu=\lecamN(0,1)$, $A(\theta)=\theta^2/2$) the remainder vanishes, and the local experiment is the Gaussian shift for every $n$. For Bernoulli and Poisson data in their natural parameters, the theorem gives $I=p(1-p)$ and $I=\lambda$.

For general models the right hypothesis is one derivative, taken in $L^2$. The model is \term{differentiable in quadratic mean}\index{differentiability in quadratic mean} at $\theta_0$ if there is a measurable function $\dot\ell_{\theta_0}$ such that, as $t\to0$,
\[
  \int\Big(\sqrt{p_{\theta_0+t}}-\sqrt{p_{\theta_0}}-\tfrac12t\,\dot\ell_{\theta_0}\sqrt{p_{\theta_0}}\Big)^2\dd\mu=o(t^2).
\]

\begin{theorem}[differentiability in quadratic mean implies LAN]\label{thm:lecam:dqm}
If the model is differentiable in quadratic mean at $\theta_0$, then $\E_{\theta_0}\dot\ell_{\theta_0}=0$ and $I(\theta_0)=\E_{\theta_0}\dot\ell_{\theta_0}^2<\infty$. If moreover $I(\theta_0)>0$, then the model is LAN at $\theta_0$, with $\Delta_n=n^{-1/2}\sum_i\dot\ell_{\theta_0}(X_i)$ and $I_0=I(\theta_0)$.
\end{theorem}

We do not reprove this \parencite[Chapter~7]{vanderVaart1998}. It gives the sense of \emph{regular} in the title of this chapter: a regular model is differentiable in quadratic mean, with positive information, at every point. The condition holds when $\theta\mapsto\sqrt{p_\theta(x)}$ is continuously differentiable for every $x$ and $I$ is finite and continuous, and also in the Laplace location model $p_\theta(x)=\frac12e^{-\abs{x-\theta}}$, where $\theta\mapsto\log p_\theta(x)$ is not differentiable at $\theta=x$ and Cramér's hypotheses fail; there $\dot\ell_\theta(x)=\sgn(x-\theta)$ and $I=1$.\sidenote{In the Laplace model the maximum likelihood estimator is the sample median, and $\sqrt n(\mathrm{med}_n-\theta)\lecamto\lecamN(0,1)=\lecamN(0,1/I)$: it is efficient although Cramér's theorem does not apply.} \textcite{LeCam1970} examined which of the classical assumptions are needed for the asymptotic normality of maximum likelihood estimates; for the theorems of Hájek in \cref{sec:lecam:transfer}, one derivative in quadratic mean is all that is needed.

\section{After: carrying the answers back}\label{sec:lecam:transfer}

\subsection{Fisher's bound, correctly stated}

The first transfer needs no general theory: the van Trees inequality holds in every model that satisfies its hypotheses, so we can run the proof of \cref{thm:lecam:gaussian-minimax} in the $n$th experiment, with the prior shrunk to the scale $1/\sqrt n$ that the limit singled out.

\begin{theorem}[local asymptotic minimax bound for squared error]\label{thm:lecam:lam-quadratic}
Suppose that the model satisfies conditions (a) and (b) of \cref{lem:lecam:vantrees} on every compact subinterval of $\Theta$, and that $I$ is continuous at $\theta_0\in\Theta$. Let $c>0$, and let $n$ be so large that $[\theta_0-c/\sqrt n,\theta_0+c/\sqrt n]\subset\Theta$. Then every estimator $T_n$ based on $X_1,\dots,X_n$ satisfies
\begin{equation}\label{eq:lecam:lam-finite}
  \sup_{\abs h\le c}\,n\E_{\theta_0+h/\sqrt n}\Big(T_n-\theta_0-\frac h{\sqrt n}\Big)^2\ \ge\ \frac1{\sup_{\abs{\theta-\theta_0}\le c/\sqrt n}I(\theta)+\pi^2/c^2}.
\end{equation}
Consequently,
\[
  \lim_{c\to\infty}\liminf_{n\to\infty}\,\inf_{T_n}\,\sup_{\abs h\le c}\,n\E_{\theta_0+h/\sqrt n}\Big(T_n-\theta_0-\frac h{\sqrt n}\Big)^2\ \ge\ \frac1{I(\theta_0)}.
\]
\end{theorem}

\begin{proof}
The densities $p^n_\theta(x_1,\dots,x_n)=\prod_ip_\theta(x_i)$ of the sample satisfy (a) and (b) with information $nI(\theta)$: the derivative of the product is $\sum_i\dot p_\theta(x_i)\prod_{j\ne i}p_\theta(x_j)$, whose integral vanishes term by term, and in its square divided by $p^n_\theta$ the cross terms integrate to $(\int\dot p_\theta\dd\mu)^2=0$. Apply \cref{lem:lecam:vantrees} to these densities on $[\theta_0-c/\sqrt n,\theta_0+c/\sqrt n]$ with the rescaled prior $\lambda(\theta)=\sqrt n\,\lambda_c(\sqrt n(\theta-\theta_0))$, where $\lambda_c$ is the prior of \cref{thm:lecam:gaussian-minimax}. A change of variables gives $J(\lambda)=nJ(\lambda_c)=n\pi^2/c^2$, and $\int nI\lambda\dd\theta$ is at most $n$ times the supremum of $I$ over the interval, so the Bayes risk of $T_n$ is at least $1/(n\sup I+n\pi^2/c^2)$. Multiplying by $n$ and bounding the supremum over $\abs h\le c$ below by the average gives \eqref{eq:lecam:lam-finite}. As $n\to\infty$ its right side tends to $1/(I(\theta_0)+\pi^2/c^2)$ by the continuity of $I$ at $\theta_0$, and this tends to $1/I(\theta_0)$ as $c\to\infty$.
\end{proof}

Compare this with the restricted-class aside of \cref{sec:lecam:before}: no assumption on the estimator, no Neyman--Pearson lemma, no expansion of a likelihood ratio. The price is the form of the statement, a worst case over a neighborhood of size $c/\sqrt n$ instead of a value at a point, and Hodges's estimator shows that it must be paid. The theorem also quantifies Le Cam's observation about superefficiency: however good an estimator is at $\theta_0$, for every $c$ and $\eps>0$ and all large $n$ its normalized risk exceeds $1/(I(\theta_0)+\pi^2/c^2)-\eps$ somewhere within $c/\sqrt n$ of $\theta_0$. The sample mean of the normal model, with normalized risk $1$ everywhere, attains the bound.

\subsection{Contiguity and the representation theorem}

For other losses, for tests, and for limit laws, the transfer goes through the experiments themselves, by a change of measure. A sequence of probability measures $Q_n$ is \term{contiguous}\index{contiguity} to $P_n$ if $P_n(A_n)\to0$ implies $Q_n(A_n)\to0$ for every sequence of events $A_n$. Contiguity and the next two lemmas are due to \textcite{LeCam1960}; \textcite[Chapter~6]{vanderVaart1998} gives proofs.

\begin{theorem}[Le Cam's first and third lemmas]\label{thm:lecam:third-lemma}\index{Le Cam's third lemma}
Let $P_n$ and $Q_n$ be probability measures on a common space, and let $L_n=\dd Q_n/\dd P_n$ be the density of the absolutely continuous part of $Q_n$ with respect to $P_n$.
\begin{enumerate}[label=(\alph*)]
\item If $\log L_n\lecamto\lecamN(-\tfrac12\tau^2,\tau^2)$ under $P_n$, then $Q_n$ is contiguous to $P_n$.
\item If real statistics $Y_n$ satisfy $(Y_n,\log L_n)\lecamto\lecamN_2\big((\mu,-\tfrac12\tau^2),\big(\begin{smallmatrix}\sigma^2&\kappa\\\kappa&\tau^2\end{smallmatrix}\big)\big)$ under $P_n$, then $Y_n\lecamto\lecamN(\mu+\kappa,\sigma^2)$ under $Q_n$.
\end{enumerate}
\end{theorem}

Under LAN take $P_n=P^n_{\theta_0}$ and $Q_n=P^n_{\theta_0+h/\sqrt n}$. By \eqref{eq:lecam:lan} the log-likelihood ratio is asymptotically $\lecamN(-\frac12h^2I_0,h^2I_0)$, so local alternatives are contiguous; and since $\Delta_n$ has asymptotic covariance $hI_0$ with it, $\Delta_n/I_0\lecamto\lecamN(h,1/I_0)$ under $\theta_0+h/\sqrt n$. The observation of the limit experiment is carried by the sequence.

\begin{theorem}[asymptotic representation]\label{thm:lecam:representation}
Assume LAN at $\theta_0$ with information $I_0$, and let $T_n$ be estimators such that for every $h$ the law of $\sqrt n(T_n-\theta_0-h/\sqrt n)$ under $P^n_{\theta_0+h/\sqrt n}$ converges to a distribution $L_h$. Then there is a randomized estimator $T$ in the Gaussian shift $(\lecamN(h,1/I_0):h\in\R)$ such that $T-h$ has the law $L_h$ under $h$, for every $h$.
\end{theorem}

\begin{proofsketch}
Under $\theta_0$ the pair $(S_n,\Delta_n)$, with $S_n=\sqrt n(T_n-\theta_0)$, is tight, so along a subsequence it converges in law to some $(S,\Delta)$ with $\Delta\sim\lecamN(0,I_0)$. For bounded continuous $f$, contiguity gives $\E_{\theta_0+h/\sqrt n}f(S_n)=\E_{\theta_0}[f(S_n)L_n(h)]+o(1)$, with $L_n(h)$ the likelihood ratio in \eqref{eq:lecam:lan}, and by LAN and uniform integrability the right side tends to $\E[f(S)e^{h\Delta-h^2I_0/2}]$; this is also how part~(b) of \cref{thm:lecam:third-lemma} is proved. The limit law of $S_n$ under $\theta_0+h/\sqrt n$ is thus the law of $S$ tilted by a density that depends only on $\Delta$. Since $I_0X$, for $X\sim\lecamN(h,1/I_0)$, has exactly this density relative to $\lecamN(0,I_0)$, drawing $T$ from the conditional law of $S$ given $\Delta=I_0X$ gives $T$ the law $L_h$ shifted by $h$. Omitted are the uniform integrability and the conditional laws \parencite[Chapters~7 and~9]{vanderVaart1998}; \textcite{LeCam1972} treats limits of experiments in general.
\end{proofsketch}

\subsection{Hájek's theorems and Le Cam's null sets}

\begin{definition}[regular estimator]\label{def:lecam:regular}
An estimator sequence $T_n$ is \term{regular}\index{regular estimator} at $\theta_0$ if there is a distribution $L$ such that $\sqrt n(T_n-\theta_0-h/\sqrt n)\lecamto L$ under $P^n_{\theta_0+h/\sqrt n}$ for every $h\in\R$.
\end{definition}

The term is Hájek's.\sidenote{The word now has two uses: a model is regular if it is differentiable in quadratic mean with positive information, and an estimator sequence is regular in Hájek's sense. Both say that nothing jumps under a perturbation of size $1/\sqrt n$.} The sample mean of the normal model is regular everywhere; Hodges's estimator is regular at every $\theta_0\ne0$ but not at $0$, where its local limit law is the point mass at $-h$.

\begin{theorem}[convolution theorem]\label{thm:lecam:convolution}\index{convolution theorem}
Assume LAN at $\theta_0$ with information $I_0$. If $T_n$ is regular at $\theta_0$ with limit law $L$, then $L=\lecamN(0,1/I_0)*M$ for some distribution $M$. In particular, if $L$ has finite variance, then that variance is at least $1/I_0$.
\end{theorem}

\begin{proofsketch}
By \cref{thm:lecam:representation} there is a randomized estimator $T$ in the Gaussian shift with $T-h\sim L$ for every $h$, and \cref{prop:lecam:convolution-limit} applies to it. Only the proof of \cref{thm:lecam:representation} is omitted.
\end{proofsketch}

This is the theorem of \textcite{Hajek1970}; \textcite{Inagaki1970} found a similar result independently under stronger assumptions. A regular estimator with limit law $\lecamN(0,1/I_0)$, such as the maximum likelihood estimator of an exponential family (Exercise~\ref{ex:lecam:mle}), is best among regular estimators in a strong sense: every other one is, in the limit, the best one plus independent noise. Regularity is the honest form of the restricted classes of \cref{sec:lecam:before}, because the limit experiment explains it as equivariance in law.

\begin{theorem}[local asymptotic minimax theorem]\label{thm:lecam:lam}\index{local asymptotic minimax theorem}
Assume LAN at $\theta_0$ with information $I_0$, and let $\ell:\R\to[0,\infty)$ be symmetric and nondecreasing on $[0,\infty)$. Then every estimator sequence $T_n$ satisfies
\[
  \lim_{c\to\infty}\liminf_{n\to\infty}\,\sup_{\abs h\le c}\,\E_{\theta_0+h/\sqrt n}\,\ell\big(\sqrt n(T_n-\theta_0-h/\sqrt n)\big)\ \ge\ \E\,\ell\big(Z/\sqrt{I_0}\big),
\]
with $Z$ standard normal.
\end{theorem}

\begin{proofsketch}
Truncate $\ell$, fix a finite set $F\subset[-c,c]$, and pass to a subsequence that attains the $\liminf$ and along which the local laws converge for $h\in F$, after compactifying the line. A version of \cref{thm:lecam:representation} for finitely many $h$ bounds the left side below by the minimax risk of the Gaussian shift on $F$, and these risks increase to $\E\ell(Z/\sqrt{I_0})$ as $F$ grows (and $c$ with it), by a Bayes argument with normal priors. The details are in \textcite{Hajek1972} and \textcite[Chapter~8]{vanderVaart1998}. For $\ell(x)=x^2$, and under its own hypotheses, \cref{thm:lecam:lam-quadratic} gives a complete, nonasymptotic proof.
\end{proofsketch}

The second line of defense of \cref{sec:lecam:before} also becomes short: a Fubini argument followed by the convolution theorem in the limit.

\begin{theorem}[superefficiency happens on a null set]\label{thm:lecam:superefficiency}
Assume that the model is differentiable in quadratic mean at every $\theta\in\Theta$, with $I(\theta)>0$. Let $T_n$ be estimators such that $\sqrt n(T_n-\theta)\lecamto L_\theta$ under $P^n_\theta$ for every $\theta\in\Theta$. Then for Lebesgue-almost every $\theta$ there is a distribution $M_\theta$ with $L_\theta=\lecamN(0,1/I(\theta))*M_\theta$; in particular the variance of $L_\theta$, when finite, is at least $1/I(\theta)$ for almost every $\theta$.
\end{theorem}

\begin{proofsketch}
Pointwise convergence forces regularity almost everywhere. For bounded continuous $f$ put $g_n(\theta)=\E_\theta f(\sqrt n(T_n-\theta))$ and $g(\theta)=\int f\dd L_\theta$, so that $g_n(\theta+h/\sqrt n)$ is the expectation of $f(\sqrt n(T_n-\theta-h/\sqrt n))$ under $\theta+h/\sqrt n$. Since $g_n\to g$ boundedly and translation is continuous in $L^1$, the integral $\int\abs{g_n(\theta+h/\sqrt n)-g(\theta)}w(\theta)\dd\theta$ tends to zero for every $h$ and every probability density $w$ with compact support in $\Theta$. Fubini's theorem, a subsequence, and a countable convergence-determining family of functions $f$ then make the sequence regular at almost every $\theta$, along a subsequence and for almost every $h$, and the representation theorem and \cref{prop:lecam:convolution-limit} finish. Omitted are measurability, the subsequences, and the removal of the exceptional $h$ \parencite[Chapter~8]{vanderVaart1998}.
\end{proofsketch}

For Hodges's estimator the exceptional set is $\{0\}$. The theorem is a modern form of Le Cam's null-set theorem: at almost every point the limit law is that of the efficient estimator plus independent noise.

\subsection{New reach}

\emph{Tests.} In $\lecamN(h,1/I_0)$ the most powerful test of $h=0$ against a fixed $h>0$ rejects for large $X$, whatever $h$ is. Transferred back with the Neyman--Pearson lemma and \cref{thm:lecam:third-lemma}, this says that no test with asymptotic level $\alpha$ has asymptotic power above $1-\Phi(z_\alpha-h\sqrt{I_0})$ at $\theta_0+h/\sqrt n$, with $z_\alpha=\Phi^{-1}(1-\alpha)$, and that the score test, rejecting when $\Delta_n/\sqrt{I_0}>z_\alpha$, attains this envelope for every $h>0$ at once.

\emph{Semiparametric efficiency.}\index{semiparametric efficiency} \textcite{Stein1956} observed that a nonparametric estimation problem is at least as hard as each of its one-dimensional parametric submodels. Localizing along every smooth path through the truth, each a LAN model, turns this into a theory of information bounds and efficient estimators \parencites{BickelKlaassenRitovWellner1993}[Chapter~25]{vanderVaart1998}.

\emph{Stochastic optimization.} \textcite{DuchiRuan2021} develop a local minimax theory for stochastic convex optimization in the mold of Hájek and Le Cam. In smooth unconstrained problems its bounds are attained by averaged stochastic gradient methods \parencite{PolyakJuditsky1992}; with constraints, dual averaging identifies the active constraints in finite time, while plain stochastic gradient methods can be suboptimal. Exercise~\ref{ex:lecam:sgd} localizes inside a Gaussian family instead and finds a smaller bound.

\section{Costs and limits}\label{sec:lecam:costs}

\paragraph{Where did the difficulty go?}
It was paid in three places, each once. The expansion: differentiability in quadratic mean implies LAN by an $L^2$ argument with a careful remainder, after which each model costs one derivative. The transfer: contiguity, the third lemma, and the representation theorem combine a change of measure, uniform integrability, and Prohorov's theorem, once for all models, losses, and problems. And, beyond LAN, Le Cam's general theory of deficiency and limits of experiments \parencite{LeCam1972,LeCam1986}. The old route paid for third derivatives and Neyman--Pearson computations in each model and for each question, so this is a bargain.

\paragraph{What did not dissolve.}
\emph{Nonregular models.} When the support moves with the parameter, the rate and the limit both change: for the uniform distribution on $[0,\theta]$ the scale is $1/n$, and the limit experiment observes $h-E$ with $E$ exponential (Exercise~\ref{ex:lecam:uniform}), a location experiment that is not Gaussian. Boundary points, change points, and vanishing information give other limits, each to be found and solved on its own \parencite{IbragimovHasminskii1981}. \emph{Many parameters.} When the dimension grows with $n$ the local Gaussian picture can break: in spiked covariance models the log-likelihood-ratio process converges below a detection threshold to a Gaussian process, but the limit experiment is not a Gaussian shift, and the experiments separate above the threshold \parencite{OnatskiMoreiraHallin2013}; for regression with a number of parameters proportional to $n$ we know of no comparable theory (Exercise~\ref{ex:lecam:proportional}). \emph{Finite samples.} Limit statements say nothing about $n=50$ unless the distance between $\cE_n$ and its limit is quantified. \emph{Local is not global.} In dimension $d\ge3$ the limit $\lecamN(h,I_d)$ shows Stein's phenomenon: the observation is minimax but inadmissible, since the James--Stein estimator\index{James--Stein estimator} has smaller risk at every $h$ \parencite{JamesStein1961}. The local asymptotic minimax theorem determines the minimax risk, not the admissible procedures, and some superefficient estimators pay nothing anywhere (Exercise~\ref{ex:lecam:stein}).

\paragraph{Other seas.}
\textcite{IbragimovHasminskii1981} built a parallel theory on likelihood-ratio processes that treats many nonregular models directly. Bayesian procedures localize in the same way: under LAN together with identifiability and prior conditions, the posterior of $h$ is asymptotically $\lecamN(\Delta_n/I_0,1/I_0)$, the Bernstein--von Mises theorem \parencite[Chapter~10]{vanderVaart1998}. Nonasymptotic minimax theory keeps the localization but replaces the limit experiment by small finite families of hypotheses, as in Le Cam's two-point method.

\section{The move, distilled}\label{sec:lecam:distilled}

\begin{move}{Localize to a limit object}
\movepart{In one sentence} When a sequence of problems has no best answer, rescale it where it neither trivializes nor degenerates, identify the limit object, solve the problem there, and carry the solution back with a transfer theorem proved once.
\movepart{Triggers} Statements about a single parameter value that are defeated by procedures tuned to that value; answers that depend on the problem only through one or two numbers, such as $I(\theta_0)$; proofs that repeat the same second-order expansion in every model.
\movepart{Recipe} (1) Find the scale at which neighboring problems stay at distance of order one. (2) Rescale and compute the limit of the likelihood ratios: the limit object. (3) Solve the problem there exactly, by Bayes, minimax, or equivariance arguments. (4) Prove once a transfer theorem from the limit to sequences. (5) Construct procedures that converge to the optimal one in the limit.
\movepart{Anti-pattern} Localizing at the wrong scale: at a fixed parameter value, which lets Hodges back in, or at a finer scale, where every estimator looks optimal. Taking limits that do not transfer, such as pointwise asymptotic variances. Declaring the limit Gaussian out of habit.
\end{move}

\dissolutiontest{Fisher's bound, in local minimax form, is one inequality with a shrinking prior (\cref{thm:lecam:lam-quadratic}) instead of two expansions per model; Hodges's paradox is one line in the limit.}{No unbiasedness, no uniformity conditions on estimators, no third derivatives; Hájek's theorems need one derivative in quadratic mean.}{Cramér--Rao, efficiency of the maximum likelihood estimator, Hájek's two theorems, and power envelopes are facts about one experiment, $\lecamN(h,1/I_0)$.}{Superefficiency confined to null sets; semiparametric efficiency; nonregular limits; local minimax theory for stochastic optimization.}

The move is a cousin of completion (\cref{ch:lebesgue,ch:graphon}): there a sequence of functions or graphs is replaced by its limit in a completed space, here a sequence of experiments by its limit in Le Cam's space of experiments, and for regular models at the scale $1/\sqrt n$ the limit is always a Gaussian shift. LAN is a central limit theorem for experiments. The lower bounds of \cref{thm:lecam:gaussian-minimax,thm:lecam:lam-quadratic} are Yao's principle (\cref{ch:lpduality}) at the scale of the noise: a prior that is nearly least favorable in the limit certifies a bound for every procedure. \Cref{ch:conformal} takes the opposite route, buying finite-sample validity from a symmetry and claiming no optimality; Le Cam's theory buys optimality and pays with asymptotics.

\begin{exercises}
\item \consolidation (LAN by hand) For Bernoulli data with success probability $p\in(0,1)$, show by a direct Taylor expansion that the model is LAN at every $p_0$, with $I_0=1/(p_0(1-p_0))$ and $\Delta_n=I_0\,n^{-1/2}\sum_i(X_i-p_0)$; do the same for Poisson$(\lambda)$ data, with $I_0=1/\lambda_0$ and $\Delta_n=I_0\,n^{-1/2}\sum_i(X_i-\lambda_0)$. Then recover the Bernoulli result from \cref{thm:lecam:lan-expfam} in the log-odds $\theta$, using the uniformity of its remainder: LAN survives the change of parameter, with information multiplied by $(\mathrm d\theta/\mathrm dp)^2$.\label{ex:lecam:bernoulli}

\item \consolidation\hard (Hodges's estimator revisited) Replace $n^{-1/4}$ in Hodges's estimator by thresholds $a_n>0$. (a) Show that if $a_n\to0$ and $\sqrt n\,a_n\to\infty$, then \cref{prop:lecam:hodges,prop:lecam:hodges-local} remain true, with $\sqrt n/3$ replaced by $na_n^2/3$ whenever $\sqrt n\,a_n\ge1$. (b) Show that if $\sqrt n\,a_n\to a\in(0,\infty)$, then under $h/\sqrt n$ the law of $\sqrt n(T_n-h/\sqrt n)$ tends to that of $X\ind\{\abs X\ge a\}-h$ with $X\sim\lecamN(h,1)$, so that $T_n$ is not regular at $0$. (c) For the original estimator, show that its bias $b_n$ satisfies $b_n'(0)\to-1$, so that $n$ times the Cramér--Rao bound at $0$ tends to zero. \hint{for (c), write $\E_\theta T_n$ as an integral against $\varphi(y-\sqrt n\,\theta)$ and differentiate under the integral sign.}\label{ex:lecam:hodges}

\item \consolidation\hard (the best prior) Show that every absolutely continuous probability density $\lambda$ on $[-c,c]$ with $\lambda(\pm c)=0$ and $J(\lambda)<\infty$ satisfies $J(\lambda)\ge\pi^2/c^2$, with equality for $\lambda_c$. \hint{put $g=\sqrt\lambda$, so that $J(\lambda)=4\int g'^2$; expand $g$ in the sine series $\sum_kb_k\sin(k\pi(h+c)/2c)$ and use Parseval's identity for $g$ and for $g'$.}\label{ex:lecam:prior}

\item \consolidation (the scale $1/\sqrt n$) Show that the Hellinger affinity $\rho(P,Q)=\int\sqrt{pq}\dd\mu$ of $\lecamN(a,1)$ and $\lecamN(b,1)$ is $e^{-(a-b)^2/8}$, that $\rho(P^n,Q^n)=\rho(P,Q)^n$, and that $1-\sqrt{2-2\rho}\le\int\min(p,q)\dd\mu\le\rho$ in general. Deduce that the smallest sum of the error probabilities of a test of $\lecamN(\theta,1)^n$ against $\lecamN(\theta_n',1)^n$ tends to zero if $\sqrt n\abs{\theta-\theta_n'}\to\infty$ and to one if $\sqrt n\abs{\theta-\theta_n'}\to0$.\label{ex:lecam:hellinger}

\item \consolidation (efficiency of the maximum likelihood estimator) In the setting of \cref{thm:lecam:lan-expfam}, let $\hat\theta_n=(A')^{-1}(\bar S_n)$ when $\bar S_n=n^{-1}\sum_iS(X_i)$ lies in the open interval $A'(\Theta)$, and let $\hat\theta_n$ be a fixed point of $\Theta$ otherwise. Show that $\sqrt n(\hat\theta_n-\theta_0)=\Delta_n/I_0+o_P(1)$ under $\theta_0$, and use \cref{thm:lecam:third-lemma} to show that $\hat\theta_n$ is regular at $\theta_0$ with limit law $\lecamN(0,1/I_0)$.\label{ex:lecam:mle}

\item \transfer (a non-Gaussian limit) Let $X_i$ be uniform on $[0,\theta]$, $\theta>0$, and localize at rate $n$: $\theta=\theta_0+h/n$. (a) Show that $n(X_{(n)}-\theta_0)\lecamto h-E$ under $\theta_0+h/n$, where $X_{(n)}$ is the sample maximum and $E$ is exponential with mean $\theta_0$. (b) Show that the density of the absolutely continuous part of $P^n_{\theta_0+h/n}$ with respect to $P^n_{\theta_0}$ converges in law under $\theta_0$ to $e^{-h/\theta_0}\ind\{-E\le h\}$, the likelihood ratio at base point $0$ of the experiment that observes $Z=h-E$. Why are the alternatives with $h>0$ not contiguous? (c) In the limit, find the best estimator of the form $Z+b$ for squared error, and compare its risk $\theta_0^2$ with the risk $2\theta_0^2$ of $Z$. Translate back to $X_{(n)}$ and $(1+1/n)X_{(n)}$.\label{ex:lecam:uniform}

\item \transfer\hard (Stein's phenomenon in the limit) Let $X\sim\lecamN(h,I_d)$ with $d\ge3$ and $\delta(X)=(1-(d-2)/\norm X^2)X$. (a) Using $\E[(X-h)\cdot g(X)]=\E[\operatorname{div}g(X)]$ for $g(x)=x/\norm x^2$, show that $\E\norm{\delta(X)-h}^2=d-(d-2)^2\E\norm X^{-2}<d$ for every $h$. (b) For independent $\lecamN(\theta,I_d)$ observations let $T_n=\theta_0+n^{-1/2}\delta(\sqrt n(\bar X_n-\theta_0))$. Show that $n\E_\theta\norm{T_n-\theta}^2$ at $\theta=\theta_0+h/\sqrt n$ is the risk of $\delta$ at $h$: it is $2$ at $h=0$, below $d$ everywhere, and its supremum over $\norm h\le c$ tends to $d$ as $c\to\infty$. (c) Show that $T_n$ is not regular at $\theta_0$, although it is superefficient there without paying anywhere. \hint{$\E\norm X^{-2}=1/(d-2)$ when $h=0$, and $\E\norm X^{-2}\to0$ as $\norm h\to\infty$.}\label{ex:lecam:stein}

\item \transfer\hard (local lower bounds for stochastic optimization) Let $\rho:\R\to\R$ be even and twice continuously differentiable, with $0<\rho''\le\beta$ for a constant $\beta$. A method observes independent $\xi_1,\dots,\xi_n\sim\lecamN(\theta,\sigma^2)$ and seeks the minimizer $x^*(\theta)$ of $f_\theta(x)=\E_\theta\rho(x-\xi_1)$. (a) Show that $x^*(\theta)=\theta$. (b) Deduce from \cref{thm:lecam:lam-quadratic} that every method, randomized or not (see the remark after \cref{lem:lecam:vantrees}), including every stochastic gradient method that queries $\rho'(x-\xi_i)$, has local asymptotic minimax risk at least $\sigma^2$ for $n(\hat x_n-x^*)^2$. (c) Averaged stochastic gradient descent with suitable step sizes has asymptotic variance $\Sigma/H^2$, with $H=\E\rho''(\sigma Z)$, $\Sigma=\E\rho'(\sigma Z)^2$, and $Z$ standard normal \parencite{PolyakJuditsky1992}. Show that $\Sigma/H^2\ge\sigma^2$, with equality only for quadratic $\rho$. What does the bound $\sigma^2$ know that gradients do not (compare \textcite{DuchiRuan2021})? \hint{Stein's identity $\E[Z\rho'(\sigma Z)]=\sigma\E\rho''(\sigma Z)$ and the Cauchy--Schwarz inequality.}\label{ex:lecam:sgd}

\item \research\harder (limit experiments in proportional asymptotics) In logistic regression and robust linear regression with $p$ parameters, $n$ observations, and $p/n\to\kappa\in(0,1)$, the classical $1/I$ theory fails. Define local experiments for such sequences, for instance by localizing one linear functional of the parameter at rate $1/\sqrt n$; decide whether they converge, to what, and whether a convolution theorem holds. In spiked covariance models a limit experiment exists that is not a Gaussian shift \parencite{OnatskiMoreiraHallin2013}; for proportional regression a general theory of this kind is, as far as we know, open. \hint{start from Gaussian linear regression, where the experiment is exactly Gaussian.}\label{ex:lecam:proportional}
\end{exercises}

\notesandrefs
Fisher's program is in \textcite{Fisher1922,Fisher1925}. \textcite{Vovk2009} reviews the early history of superefficiency: Hotelling's letter of 1930, Hodges's examples of 1951, Le Cam's announcement of the null-set theorem (a talk at a meeting of the Institute of Mathematical Statistics in December 1951), the flaw in the thesis proof, and the three lines of defense; \textcite{vanderVaart1997} gives several proofs of the theorem.

Local asymptotic normality and contiguity come from \textcite{LeCam1960}. \textcite{Hajek1970} proved the convolution theorem under LAN, and \textcite{Inagaki1970} a similar result independently. Hájek's local asymptotic minimax theorem \parencite{Hajek1972} and Le Cam's theory of limits of experiments \parencite{LeCam1972} appeared in the proceedings of the Sixth Berkeley Symposium, held in 1970. \textcite{LeCam1986} is the comprehensive account, \textcite{LeCamYang2000} a shorter introduction, and \textcite{vanderVaart1998} the standard textbook; its Chapters~6 to~9 are the source of our statements in \cref{sec:lecam:transfer}.

The van Trees inequality is from \textcite{vanTrees1968}, and \textcite{GillLevit1995} developed its uses in asymptotic statistics; \cref{thm:lecam:lam-quadratic} is in their spirit. Semiparametric efficiency goes back to \textcite{Stein1956}, with \textcite{BickelKlaassenRitovWellner1993} as the standard reference.

\printchapterbib
\end{refsection}
